\answer{ex:03:1}{$g_E=g_L$で$35\to17.5\mV$、$g_E=4g_L$で$56\to40\mV$。引き算なら$38.5\mV$になる。シャントは$g_E$を$3$で割る。窓は半分に狭まる（$\tau_{\text{eff}}$が$2$分の1になる）。}
\answer{ex:03:2}{$\operatorname{tr}=\frac{w_{EE}-1}{\tau_E}-\frac1{\tau_I}$、$\det=\frac{1-w_{EE}+w_{EI}w_{IE}}{\tau_E\tau_I}$。境界では$\omega=\sqrt{\det}$、$\tau_I=20\ms$、周期$73\ms$。}
\answer{ex:03:3}{$d_c=\tau_E\arccos(-1/G)/\sqrt{G^2-1}$、周期$2\pi\tau_E/\sqrt{G^2-1}\in(2d_c,4d_c)$。$G=2$：$d_c=12.1\ms$、周期$36\ms$。$G=5$：$d_c=3.6\ms$、周期$12.8\ms$で、ちょうど速いリズムの帯域である。}
\answer{ex:03:4}{(a)~上げるときは$I_2=\beta I_1=2$、下げるときは$I_2=I_1/\beta=0.5$で、幅$1.5$のヒステリシス。(b)~$\varepsilon$が1桁小さくなるごとに$\tau\ln10/(\beta-1)=2.3\tau$ずつ増える。}
\answer{ex:03:5}{3個、$d_k=5$、$10$、$15\ms$。禁止は$15\ms$を覆う必要があり、1個あたり$5\ms$。}
\answer{ex:03:6}{$n+M+1=3+4+1=8$個。2個なら：\nt{GABA} $=[x+y+z\ge2]$、出力 $=[x+y+z-2\,\nt{GABA}\ge1]$。}
\answer{ex:03:7}{$\binom84=70<100\le\binom94=126$より仲介役は$9$個（シュペルナーの定理）。直接の結合がなければ$100$個：$\beta(J-I)$の階数が$n$だから。}
\answer{ex:03:8}{$h_c=\frac1{2w_{EE}}+\frac{w_{EI}w_{IE}-w_{EE}}{4w_{EE}^2}=\frac7{18}\approx0.39$。シミュレーションでは$h=0.38$と$0.40$の間で符号が変わる。}
